\section{Finite random projection and retention of the original shadings}\label{sec:projection}

\begin{lemma}[Directional collisions and nonshort projections]\label{lem:projection}
Fix $M\ge4$, $s\ge3$, and a constant $a>0$.
If $m$ original $\delta$-separated directions satisfy\eqref{eq:angular-cap}, then there exist
$P:\R^M\to\R^4$ with $\norm P\le C_M$ and a subfamily $\T_*$ of the original labels such that
\begin{align}
 |Pv_T|&\ge c_M &&(T\in\T_*),\label{eq:nonshort}\\
 \ang(Pv_T,Pv_{T'})&>a\delta &&(T\ne T'\in\T_*),\label{eq:projected-separated}\\
 \card{\T_*}&\ge
 \frac{c_{M,s,a}m}{A\delta^{3-s}\Ldelta}.\label{eq:projection-count}
\end{align}
\end{lemma}

\begin{proof}
Let the entries of $P$ be independent standard real Gaussian variables.
Choose $B_M$ sufficiently large that $\Pr\{\norm P>B_M\}<1/32$.
For a fixed unit vector $u$, $Pu$ is a standard four-dimensional Gaussian vector.
We can therefore choose $c_M>0$ sufficiently small that $\Pr\{|Pu|<c_M\}<1/128$.

If $u,v$ have projective angle $\theta$, then, after changing the sign of $v$ if necessary, we can write
\[
 Pv=(\cos\theta)Pu+(\sin\theta)H,
\]
where $H$ is a standard four-dimensional Gaussian vector independent of $Pu$.
Conditionally on $Pu$, the projection of $H$ onto $(Pu)^\perp$ is a standard three-dimensional Gaussian.
On the event $\norm P\le B_M$, a projected directional angle at most $a\delta$ implies
\[
 \big|\operatorname{proj}_{(Pu)^\perp}H\big|
 \le \frac{C_{M,a}\delta}{\sin\theta}.
\]
The three-dimensional Gaussian small-ball bound consequently gives
\begin{equation}\label{eq:collision-probability}
 \Pr\{\norm P\le B_M,\ \ang(Pu,Pv)\le a\delta\}
 \le C_{M,a}\min\{1,(\delta/\theta)^3\}.
\end{equation}
The norm constraint has not been treated as an independent event; the event $Pu=0$ has probability zero.

For each original direction, decompose into angular annuli $\theta\asymp2^j\delta$.
By\eqref{eq:angular-cap}, the expected number of collisions is at most
\begin{equation}\label{eq:collision-degree}
 C_{M,s,a}A\sum_{0\le j\le C+\log_2(1/\delta)}2^{j(s-3)}
 \le Q:=C_{M,s,a}A\delta^{3-s}\Ldelta.
\end{equation}
The contribution of projective angles larger than a fixed constant is controlled in the same way,
using a fixed finite cover of the direction space and $m\le C_MA\delta^{-s}$.
The logarithmic loss is present when $s=3$; for $s>3$ the slightly weaker uniform expression above suffices.

The expected number of short directions is less than $m/128$.
The expected number of collision-graph edges counted on the bounded-norm event is at most $CmQ$.
Applying Markov's inequality to these counts and combining the failure probabilities,
we can choose a single $P$ with bounded norm, at least $3m/4$ nonshort directions,
and at most $CmQ$ collision-graph edges among these directions.
The graph vertices are original labels, with an edge whenever the projected angle is at most $a\delta$.
The independent-set bound
\[
 \alpha(G)\ge\frac{|V(G)|^2}{|V(G)|+2|E(G)|}
\]
yields\eqref{eq:projection-count}, together with\eqref{eq:nonshort} and\eqref{eq:projected-separated}.
\end{proof}

\begin{remark}
The lemma does not assume that the projected directional distribution automatically inherits a Frostman upper bound.
Only after selecting an independent set in the collision graph do the separated directions in $\R^4$
have angular upper bounds at all scales, by volume packing in three-dimensional projective direction space.
The lemma also supplies no additional dimension along the projection kernel.
\end{remark}

\begin{lemma}[Average density, a common chart, and the visited union]\label{lem:marks}
Suppose that the original directions are $\delta$-separated and satisfy\eqref{eq:angular-cap},
and that\eqref{eq:single-line-upper} and\eqref{eq:mean-density} hold.
One can first select at least $c_{M,R}\lambda m$ original lines and then apply Lemma~\ref{lem:projection}
to obtain a subfamily $\T'$ in a common bounded four-dimensional graph chart.
On a grid of size $\rho\asymp\delta$, it satisfies
\begin{align}
 m':=\card{\T'}&\ge
 \frac{c_{M,R,s}\lambda m}{A\delta^{3-s}\Ldelta},\label{eq:heavy-output-count}\\
 \card{J'_T}&\ge c_{M,R}\lambda\delta^{-1}
       &&(T\in\T'),\label{eq:individual-output-marks}\\
 \card{E'}&\le C_{M,R,s}\card{E}.\label{eq:output-union}
\end{align}
Every output mark comes from its own original label; extending supporting lines adds no shading.
\end{lemma}

\begin{proof}
First remove the light lines with $\card{Y(T)}<(\lambda/2)\delta^{-1}$.
Their total mass is at most $\lambda m\delta^{-1}/2\le S/2$.
Together with\eqref{eq:single-line-upper}, this shows that at least $c_{M,R}\lambda m$ heavy lines remain.
Each retained line has at least $(\lambda/2)\delta^{-1}$ of its own original marks.
The directional upper bound\eqref{eq:angular-cap} persists on subfamilies without renormalization.
Apply Lemma~\ref{lem:projection} to the heavy family.

Move each original shaded cell to the nearest parameter point on its original line,
and remove a fixed multiplicity in the original line parameter.
Since $|Pv_T|\ge c_M$, every projected line has a directional coordinate whose absolute value is bounded below by a fixed positive constant.
Partition into finitely many charts according to a largest component and its sign, and choose a chart containing the most lines.
Use that coordinate as the common time variable. The other slopes and intercepts then have fixed bounds.
Only a fixed number of originally $\delta$-separated parameter marks can fall in a new time cell of length $\rho=C\delta$.
Thus each line retains the number of marks in\eqref{eq:individual-output-marks}.

All marks lie in a fixed bounded region. Divide the new time axis into finitely many unit windows,
let each line choose the window containing the most of its own marks,
and then pigeonhole the lines by window.
The number of marks per line and the number of lines each lose only a fixed proportion.
Translate the common window to $[0,1]$ and extend supporting lines if necessary, retaining only the existing marks within the window.
This also handles shadings concentrated near segment endpoints.
Choosing $a$ in the projection lemma to be a sufficiently large fixed constant ensures that the final slopes are separated at scale $\rho$.

Finally, the bounded-norm map $P$ sends each original $\delta$-cell into a set of diameter $O_M(\delta)$,
which is covered by a fixed number of four-dimensional $\rho$-cells.
Moving to nearest line points, changing charts, and regrouping the time grid increase the cell count only by fixed factors.
This gives\eqref{eq:output-union} and\eqref{eq:heavy-output-count}.
Marks from different original lines have not been spliced together on a new line.
\end{proof}
